% =====================================================================
%  lang-en.tex — the English edition.
%
%  English is the MASTER language of this thesis. content/en/*.tex is
%  where text is written; content/el/*.tex follows it. See docs/SYNC.md.
% =====================================================================
\setmainlanguage{english}
\setotherlanguage[numerals=arabic]{greek}

% \gr{...} is the mirror of \en{...} in the Greek edition: a Greek word
% inside English text, hyphenated by Greek rules.
\newcommand{\gr}[1]{\textgreek{#1}}

% In the English edition an "English term" macro is a no-op, but it must
% still exist: the two editions share figure sources and table bodies,
% and those use \en{} freely.
\newcommand{\en}[1]{\emph{#1}}
\newcommand{\enup}[1]{#1}

\newcommand{\dgblankleaftext}{(This page is intentionally blank)}

% Names. Called explicitly from main-en.tex after \begin{document},
% because polyglossia rewrites them inside \AtBeginDocument and would
% otherwise have the last word.
\newcommand{\thesisnames}{%
  \renewcommand{\chaptername}{CHAPTER}%
  \renewcommand{\appendixname}{APPENDIX}%
  \renewcommand{\figurename}{Figure}%
  \renewcommand{\tablename}{Table}%
  \renewcommand{\bibname}{REFERENCES}%
  \renewcommand{\contentsname}{TABLE OF CONTENTS}%
  \renewcommand{\listfigurename}{LIST OF FIGURES}%
  \renewcommand{\listtablename}{LIST OF TABLES}%
}
